% §5.2 The Bloom Filter From Scratch
%
% Monograph register: construction, structural FNR=0, FPR derivation with the
% bit-independence approximation (the ex:ch5-7 counter-example folded into a
% remark), and the optimal-k theorem with a self-contained proof (the ex:ch5-5
% uniqueness argument folded into the proof body). Cross-refs sec:from-atom-to-set,
% def:bernoulli-bool, sec:predicate-algebra, sec:one-sided-two-sided.
% DEFINED: sec:bloom-from-scratch, def:bloom-filter, alg:bloom-insert-query,
% thm:bloom-optimal-k, eq:bloom-fpr, eq:bloom-optimal-k, eq:bloom-optimal-fpr.

\section{The Bloom Filter From Scratch}
\label{sec:bloom-from-scratch}

\Cref{sec:from-atom-to-set} defined the approximate set abstractly and
said nothing about building one in finite space. The Bloom filter is the
canonical construction: a bit array plus $k$ hash functions, with a
false-negative rate that is structurally zero, a false-positive rate in
closed form, and a design knob whose optimal setting we derive. Every
later chapter of Part~II returns to it.

\subsection*{The construction}

Fix a universe $\Universe$, a bit count $m$, and a hash count $k$. The
structure is a bit array $b[1\,..\,m]$, initialized to zero, with $k$ hash
functions $h_1, \ldots, h_k : \Universe \to \{1, \ldots, m\}$. The
analysis treats each $h_i$ as a uniform random function and the $k$ of
them as mutually independent; real implementations approximate this with
families such as double hashing over a single digest, which
\textcite{kirsch-mitzenmacher-2008} show leaves the asymptotic
false-positive rate unchanged, so the divergence is ignored here.

\begin{definition}[Bloom filter]
\label{def:bloom-filter}
The \emph{Bloom filter}\index{Bloom filter} of parameters $(m, k)$ is the structure
$(b, h_1, \ldots, h_k)$ with two operations:
\begin{itemize}
  \item \textsc{Insert}$(x)$: for $i = 1, \ldots, k$, set $b[h_i(x)] \gets 1$.
  \item \textsc{Query}$(x)$: return $\top$ if $b[h_i(x)] = 1$ for every
    $i$; return $\bot$ otherwise.
\end{itemize}
After inserting the $n$ elements of $A \subseteq \Universe$, the predicate
$x \mapsto \textsc{Query}(x)$ is the membership report of $\ASet{A}$.
\end{definition}

Both operations run in $O(k)$ time, independent of $|A|$ and $|\Universe|$,
and the array is the only state. We follow the literature's $(m, k, n)$
convention, with $n = |A|$ the item count; where a later chapter risks
confusing this with a workload size we write $n_{\text{items}}$.

\begin{algorithm}[ht]
  \caption{Bloom filter insert and query}
  \label{alg:bloom-insert-query}
  \SetKwInOut{Input}{input}
  \SetKwInOut{Output}{output}
  \SetKwProg{func}{procedure}{}{}
  \SetKwProg{fn}{function}{}{}

  \func{\textsc{Insert}($x$)}{
    \For{$i \gets 1$ \KwTo $k$}{
      $b[h_i(x)] \gets 1$\;
    }
  }
  \BlankLine
  \fn{\textsc{Query}($x$)}{
    \For{$i \gets 1$ \KwTo $k$}{
      \lIf{$b[h_i(x)] = 0$}{\Return $\bot$}
    }
    \Return $\top$\;
  }
\end{algorithm}

\subsection*{The false-negative rate is structurally zero}

If $x \in A$ was inserted, the bits $b[h_1(x)], \ldots, b[h_k(x)]$ were set
to $1$ at insertion, and no operation ever resets a bit. Every later
\textsc{Query}$(x)$ therefore reads $k$ ones and returns $\top$. The
argument is structural: it uses only the monotonicity of the array, not
the hash family, the load, or the workload. So
\[
  \fnrate \;=\; \Prob{\textsc{Query}(x) = \bot \Given x \in A} \;=\; 0
\]
for every parameter setting. A Bloom filter is \emph{one-sided}\index{one-sided error}
(\Cref{sec:one-sided-two-sided}): one of its two rates is driven to zero by
construction rather than by tuning. The remaining question is the size of
the other.

\subsection*{The false-positive rate}

A query for $x \notin A$ returns $\top$ iff all $k$ bits $b[h_i(x)]$ are
$1$. Fix a position $j$. Each of the $kn$ hash values produced during
insertion misses $j$ with probability $1 - 1/m$, so
\[
  \Prob{b[j] = 0} \;=\; \left(1 - \tfrac{1}{m}\right)^{kn} \;\approx\; e^{-kn/m},
\]
using $\ln(1 - 1/m) \approx -1/m$, with error of order $1/m^2$ per
insertion. Write $p = 1 - e^{-kn/m}$ for the resulting bit-occupancy
probability; the load $kn/m$ is the average number of hash hits per bit. A
query reads $k$ bits, each $1$ with marginal probability $p$. Treating the
$k$ reads as independent gives the standard approximation, the working
formula of the literature\index{Bloom filter!false-positive rate},
\begin{equation}
  \fprate \;\approx\; \left(1 - e^{-kn/m}\right)^k.
  \label{eq:bloom-fpr}
\end{equation}

\begin{remark}[The bit-independence approximation]
\label{rem:bloom-independence}
The $k$ queried bits lie in the same array populated by the same $n$
insertions, so their occupancy events are not independent, as
\textcite{bose-et-al-2008} pointed out (an exact expression is given by
\textcite{christensen-roginsky-jimeno-2010}): a collision
$h_i(x) = h_{i'}(x)$ between two query hashes makes two of the reads
literally the same event, and conditioning on one queried bit being set
raises the posterior on the others. The dependence is positive and small,
so \eqref{eq:bloom-fpr} slightly \emph{under}-estimates the true rate. The
smallest witness takes $m = 4$, $n = 1$, $k = 2$: the inserted element sets
one bit with probability $\tfrac14$ (when its two hashes collide) and two
bits otherwise, and averaging $\bigl(\tfrac{|S|}{4}\bigr)^2$ over those
cases gives an exact false-positive rate of
$\tfrac14 \cdot \tfrac{1}{16} + \tfrac34 \cdot \tfrac14 = \tfrac{13}{64}
\approx 0.203$, against $(1 - e^{-1/2})^2 \approx 0.155$ from
\eqref{eq:bloom-fpr}. Both the gap and the collision probability vanish as
$m$ grows; at the practical load $kn/m \approx \ln 2$ derived below the
deviation is well under a percent of the rate, and we use
\eqref{eq:bloom-fpr} throughout. \Cref{fig:bloom-channel} summarizes the
filter as the channel it induces\index{Bloom filter!as a channel}.
\end{remark}

\begin{figure}[ht]
\centering
\begin{tikzpicture}[font=\small, >={Stealth[length=2mm]},
    cell/.style={draw, minimum width=0.5cm, minimum height=0.5cm, inner sep=0pt},
    st/.style={draw, rounded corners, minimum width=2.4cm, minimum height=0.85cm,
      align=center}]
  % Mechanism: a bit array and k = 3 probes of one key.
  \foreach \i in {1,3,4,6,8,9,11} {\node[cell, fill=black!15] (c\i) at (0.5*\i, 0) {$1$};}
  \foreach \i in {2,5,7,10,12} {\node[cell] (c\i) at (0.5*\i, 0) {$0$};}
  \node (x) at (3.0, 1.5) {key $x$};
  \draw[->] (x.south) -- node[left, pos=0.55] {$h_1$} (c3.north);
  \draw[->] (x.south) -- node[left, pos=0.55, xshift=1pt] {$h_2$} (c6.north);
  \draw[->] (x.south) -- node[right, pos=0.55] {$h_3$} (c9.north);
  \node[align=left, anchor=west] at (6.6, 0)
    {report $\top$ iff all\\ $k$ probed bits are $1$};
  % The induced binary channel.
  \begin{scope}[yshift=-1.9cm]
    \node[st] (t1) at (0.6, 0) {truth $\top$\\($x \in A$)};
    \node[st] (t0) at (0.6, -2.6) {truth $\bot$\\($x \notin A$)};
    \node[st] (r1) at (7.6, 0) {report $\top$};
    \node[st] (r0) at (7.6, -2.6) {report $\bot$};
    \draw[->, thick] (t1) -- node[above] {$1 - \fnrate = 1$} (r1);
    \draw[->, thick] (t0) -- node[below] {$1 - \fprate$} (r0);
    \draw[->, dashed] (t1.south east) -- node[pos=0.18, above, sloped] {$\fnrate = 0$} (r0.north west);
    \draw[->] (t0.north east) -- node[pos=0.18, below, sloped] {$\fprate$} (r1.south west);
    \node at (4.1, -3.75) {$\fprate \approx (1 - e^{-kn/m})^k$, \quad $\fnrate = 0$};
  \end{scope}
\end{tikzpicture}
\caption[The Bloom filter as a binary channel]{The Bloom filter as a binary channel.
Top: a query for key $x$ reads its $k$ probed bits (here $k = 3$) and reports
$\top$ iff all are set. Bottom: the induced channel from the truth of
$x \in A$ to the report. A member's probed bits were all set at insertion, so
the false-negative rate is structurally zero; a non-member is accepted only
when every probe lands on a set bit, with probability given approximately by
\eqref{eq:bloom-fpr} (\Cref{rem:bloom-independence}).}
\label{fig:bloom-channel}
\end{figure}

\subsection*{Optimal $k$ for fixed $m$ and $n$}

The bit count $m$ is the space budget and $n$ is fixed by the application;
the hash count $k$ is the design knob. Let $f(k) = (1 - e^{-kn/m})^k$,
minimized over $k > 0$ as a continuous variable. Taking logarithms and
differentiating,
\[
  \frac{d}{dk}\ln f(k)
  \;=\; \ln(1 - e^{-kn/m})
  \;+\; \frac{kn}{m}\cdot\frac{e^{-kn/m}}{1 - e^{-kn/m}}.
\]
Writing $u = e^{-kn/m}$, so that $kn/m = -\ln u$, the first-order condition
$\ln(1-u) - (\ln u)\,\tfrac{u}{1-u} = 0$ becomes, after multiplying by
$(1-u)$,
\begin{equation}
  (1 - u)\ln(1 - u) \;=\; u \ln u.
  \label{eq:bloom-foc}
\end{equation}
By inspection $u = \tfrac12$ solves \eqref{eq:bloom-foc}, both sides equal
$-\tfrac12\ln 2$. It is the only interior solution. Set
$\phi(u) = (1-u)\ln(1-u) - u\ln u$ on $(0,1)$; it is antisymmetric about
$\tfrac12$, $\phi(1-u) = -\phi(u)$, so $\phi(\tfrac12) = 0$ and $\phi \to 0$
at both endpoints. Its derivative $\phi'(u) = -\ln\!\big(u(1-u)\big) - 2$ is
symmetric about $\tfrac12$, diverges to $+\infty$ at the endpoints, and dips
to $\phi'(\tfrac12) = 2\ln 2 - 2 < 0$, so $\phi'$ changes sign exactly
twice. Hence $\phi > 0$ on $(0, \tfrac12)$ and $\phi < 0$ on $(\tfrac12, 1)$,
vanishing only at $u = \tfrac12$. The boundary behavior pins this critical
point as the minimum: $f \in (0,1)$ in the interior and $f \to 1$ at both
endpoints (as $k \to 0^+$, $\ln f \approx k\ln(kn/m) \to 0$; as
$k \to \infty$, $\ln f \approx -k\,e^{-kn/m} \to 0$), and the
second-derivative check at $u = \tfrac12$ gives
$\tfrac{d^2 \ln f}{dk^2} = 2(n/m)(1 - \ln 2) > 0$. Substituting
$u = \tfrac12$ into $kn/m = -\ln u$,
\begin{equation}
  k^* \;=\; \frac{m}{n}\ln 2,
  \label{eq:bloom-optimal-k}
\end{equation}
at which each bit is occupied with probability exactly $\tfrac12$ and
\begin{equation}
  \fprate^* \;\approx\; \left(\tfrac12\right)^{k^*}
  \;=\; 2^{-(m/n)\ln 2}
  \;=\; (0.6185\ldots)^{m/n}.
  \label{eq:bloom-optimal-fpr}
\end{equation}

\begin{theorem}[Optimal $k$ for the Bloom filter]
\label{thm:bloom-optimal-k}
\index{Bloom filter!optimal hash count}For fixed $m$ and $n$, the false-positive rate \eqref{eq:bloom-fpr} is
minimized at $k^* = (m/n)\ln 2$, where the array is half-occupied in
expectation and $\fprate^* \approx (1/2)^{k^*} = 2^{-(m/n)\ln 2}$.
\end{theorem}

The bit-per-item ratio $m/n$ is the right measure of a Bloom filter's space
budget: at the optimal $k$ the rate decays exponentially in $m/n$ with base
$0.6185$, so each additional bit per item divides the rate by about $1.62$,
and halving the rate costs $\log_2 e \approx 1.44$ bits per item.
\Cref{ch:information-bounds} compares this to the information-theoretic
floor and finds a fixed gap of $\log_2 e \approx 1.44$; the half-full array
at the optimum is the striking signature of the construction, most accurate
not when sparse but when exactly half its bits are set.

\subsection*{A worked example}

Take $m = 10000$, $n = 1000$, a comfortable $m/n = 10$ bits per item. Then
$k^* = 10\ln 2 \approx 6.93$, rounded to $k = 7$, giving
$\fprate^* \approx 0.6185^{10} \approx 0.0082$, about one false positive
per $120$ negative queries. Evaluating \eqref{eq:bloom-fpr} around the
optimum confirms the rounding and shows how shallow the minimum is:
\begin{center}
\begin{tabular}{ccc}
\hline
$k$ & $1 - e^{-kn/m}$ & $\fprate \approx (1 - e^{-kn/m})^k$ \\
\hline
$5$ & $0.3935$ & $0.00943$ \\
$6$ & $0.4512$ & $0.00844$ \\
$7$ & $0.5034$ & $0.00819$ \\
$8$ & $0.5507$ & $0.00846$ \\
$9$ & $0.5934$ & $0.00913$ \\
\hline
\end{tabular}
\end{center}
The rates at $k = 6, 7, 8$ agree to the third significant figure, and the
occupancy at $k = 7$ is $0.5034$, half-full as predicted. The minimum sits
where two pressures balance: raising $k$ widens the conjunction the query
must satisfy (pushing the rate down) but fills the array faster (pushing it
up), and the two cancel at half occupancy. Off the optimum the penalty is
real, $k = 1$ at $m/n = 10$ gives $\fprate \approx 0.095$, twelve times the
optimum, and the rate depends on $m$ and $n$ only through $m/n$ once $k$ is
optimized.

The Bloom filter is one row of the general predicate algebra:
$\fnrate = 0$ structurally, $\fprate$ in closed form, the same
$\Bernoulli{\BoolType}$ channel and chapter-4 axioms as any other
approximate set. \Cref{sec:predicate-algebra} develops that algebra in
generality; \Cref{sec:one-sided-two-sided} takes up the one-sided property
$\fnrate = 0$ this construction enjoys.
